\documentclass[11pt]{article}
\usepackage[a4paper,margin=25mm]{geometry}
\usepackage{amsmath,amssymb,amsthm,booktabs,hyperref,xurl}
\setlength{\emergencystretch}{2em}
\newtheorem{theorem}{Theorem}
\newtheorem{lemma}[theorem]{Lemma}
\newtheorem{corollary}[theorem]{Corollary}
\newcommand{\Z}{\mathbb Z}\newcommand{\Q}{\mathbb Q}
\title{A Pell half-trace square absorbs a nonzero guard}
\author{\(\mathbb Z\mathbb Q\) Diophantine Project}
\date{23 November 2026}
\begin{document}
\maketitle

\begin{abstract}
In the corrected signed binomial source, one square equation chooses
the sign of the first Pell trace, while a separate reciprocal ensures
that its parameter is nonzero. A different square equation performs
both jobs. Its right side is four times an even-index Pell half-trace
square on every intended witness, and equals two on the unwanted
zero-parameter branch. An existing product coordinate makes the new
equation quadratic. This saves one variable and one residual in each
complete table, giving general, specified, and ordinary rational-\(abc\)
quartics with 982, 1147, and 733 rational variables. The local result
is unconditional; each undecidability transfer retains all arithmetic
premises of its integer interface and its cited classical coding theorem.
\end{abstract}

\section{The replacement square}

For an integer \(a\), use the polynomial Pell sequences
\[
 X_n(a)+Y_n(a)\sqrt{a^2-1}
       =(a+\sqrt{a^2-1})^n
\]
or equivalently their usual integer polynomial recurrence. The
addition identities remain polynomial identities for every signed
parameter \(a\).

\begin{lemma}[Sign and completeness]\label{branch}
Let \(a,c,d\in\Z\) satisfy
\[
 |a|\ge2,\qquad |c|>1,\qquad d^2-(a^2-1)c^2=1.
\]
If \(\zeta\in\Q\) satisfies
\begin{equation}\label{newroot}
                 \zeta^2=2\bigl(ad+(a^2-1)c+1\bigr),
\end{equation}
then \(ad>0\). Conversely, if
\(c=Y_B(a)\), \(d=X_B(a)\), and \(B\ge1\) is odd, then
\eqref{newroot} has a rational, indeed integral, solution.
\end{lemma}
\begin{proof}
Put \(A=|a|\) and \(D=a^2-1\ge3\). The norm forces
\(d\ne0\). If \(ad<0\), then \(ad=-A|d|\). For \(c<0\),
\(Dc+1<0\), so the right side of \eqref{newroot} is negative.
For \(c>0\), the hypothesis gives \(c\ge2\), and
\[
 A^2d^2-(Dc+1)^2
   =(D+1)(Dc^2+1)-(Dc+1)^2
   =D(c-1)^2>0.
\]
Both \(A|d|\) and \(Dc+1\) are positive, hence
\(A|d|>Dc+1\). Again the right side of \eqref{newroot}
is negative. A rational square cannot be negative, proving \(ad>0\).

For the converse the addition and doubling identities give
\[
 ad+(a^2-1)c=X_{B+1}(a),\qquad
 X_{B+1}(a)+1=2X_{(B+1)/2}(a)^2.
\]
Thus \(\zeta=2X_{(B+1)/2}(a)\) satisfies \eqref{newroot}.
\end{proof}

\section{Excluding the zero-parameter branch first}

The retained corrected source has integral \(r,s,w,c,h\), rational
\(\tau\), and
\[
 Y=sM,\quad a=(wM+1)rY,\quad p=2wr^2Y^2M,\quad
 k=r+1+h(p-1),\qquad M=n^2\ge16.
\]
It also has the two relations
\begin{equation}\label{small}
 \tau^2=(p^2-1)k^2+1,
 \qquad k^2-4(c-kY)^2>0.
\end{equation}
The strict inequality is represented by the existing positivity row.
The source-prefix reconstruction establishes the integrality just
listed before any use of \(s\ne0\) or of Pell classification.

\begin{lemma}[The new square excludes \(s=0\)]\label{nonzero}
With the retained source relations above, the first Pell norm
\(d^2-(a^2-1)c^2=1\), and \eqref{newroot}, one has \(s\ne0\).
\end{lemma}
\begin{proof}
Suppose \(s=0\). Then \(Y=a=p=0\), while
\(k=r+1-h\) is an integer. The strict inequality in \eqref{small}
implies \(k\ne0\). The norm for \(\tau\) says
\(\tau^2+k^2=1\), so \(k=\pm1\), \(\tau=0\).
The inequality now gives \(1-4c^2>0\). Since \(c\) is an
integer, \(c=0\). The first Pell norm then gives \(d=\pm1\),
and \eqref{newroot} reduces to \(\zeta^2=2\), impossible over
\(\Q\). This contradiction proves the claim. It does not use the
reciprocal for the later integer witness \(i\), any sign of \(w\),
or the third Pell equation.
\end{proof}

\begin{theorem}[A guard-free equivalent source]\label{source}
In the eleven-integer-input source of \cite{CRT}, replace the private
square equation
\[
                   \zeta^2=(a-1)(d-1)
\]
by \eqref{newroot}, and delete the reciprocal equation
\(s s_{\rm inv}=1\) and its private coordinate \(s_{\rm inv}\).
The new and old systems have the same projection on their coding
prefix and instance. All other source coordinates and integer input
predicates remain as before; the private square witness can be
rechosen.
\end{theorem}
\begin{proof}
For a new rational solution, Lemma~\ref{nonzero} first recovers
\(s\ne0\). The proof of the signed binomial theorem in
\cite{Signed} then excludes \(w=0\), proves \(|a|\ge2\) and
\(|c|>B=2r+1\ge3\), and establishes the noncircular reciprocal
guard for \(i\). These arguments precede the selection of the first
Pell trace's sign. The one-quotient index theorem of \cite{CRT}
gives
\[
                   |c|=Y_B(|a|),\qquad |d|=X_B(|a|).
\]
Lemma~\ref{branch} supplies \(ad>0\), hence \(d=X_B(a)\)
because \(B\) is odd. The remainder of the signed binomial proof
therefore applies verbatim and gives \(w,s,c>0\) and the required
central-binomial divisibility and power-of-two relation. The old
reciprocal can be restored as \(s_{\rm inv}=1/s\). The old square
witness can also be restored: for \(B=2r+1\), the polynomial
identity
\[
 (a-1)(X_B(a)-1)
   =\bigl((a-1)(Y_{r+1}(a)+Y_r(a))\bigr)^2
\]
gives it explicitly. Since that old square coordinate occurred only
in its single defining row, rechoosing it affects nothing else.

Conversely, every old source solution reconstructs the positive
canonical first Pell values \(c=Y_B(a)\), \(d=X_B(a)\).
Lemma~\ref{branch} gives the new square witness. Drop
\(s_{\rm inv}\), keep every other coordinate, and every new
relation holds. No integer input is added or dropped in this step.
\end{proof}

\section{Quadratic implementation and complete counts}

Every predecessor table already contains the register
\(A_c=ac\) with the row \(A_c-ac=0\). Equation
\eqref{newroot} is therefore implemented by the single quadratic
residual
\begin{equation}\label{quad}
                    \zeta^2-2ad-2aA_c+2c-2.
\end{equation}
It replaces one old quadratic residual and uses the same private
square coordinate. Removing \(s_{\rm inv}\) and its row saves
one coordinate and one row. Exact occurrence checks certify that
\(s_{\rm inv}\) occurred only in the deleted guard, and that the
square coordinate occurs only in the changed row.

Applying this transformation to the complete v95 tables gives
\begin{center}
\begin{tabular}{lrrrr}
\toprule
Interface&Old variables&New variables&Old rows&New rows\\
\midrule
General&983&982&855&854\\
Specified&1148&1147&998&997\\
Ordinary rational \(abc\)&734&733&643&642\\
\bottomrule
\end{tabular}
\end{center}
The input \(x\) remains a coefficient; \(z,u,y\) remain the fixed
coding parameters. Coefficients remain in \(\Z[x,z,u,y]\), and
the sum of the squared residuals still has degree exactly four.
Theorem~\ref{source} and the previously proved complete transfer
give undecidability under exactly the arithmetic premises of each
retained integer interface and the same cited Jones coding theorem.
The local guard absorption itself is unconditional.

\section{Reproducibility}

Run
\begin{verbatim}
python undecidability/pell_halftrace_v96/build_halftrace.py
python undecidability/pell_halftrace_v96/verify_halftrace.py
\end{verbatim}
The independent verifier checks both polynomial identities, signed
Pell fixtures, the zero-\(s\) obstruction, all private occurrences,
every unchanged residual, complete coordinate support, quadratic
degree, counts, hashes, and deterministic rebuilding. The proof of
the quantified source equivalence is Theorem~\ref{source}; finite
fixtures do not replace that proof.

\begin{thebibliography}{9}
\small\setlength{\itemsep}{2pt}
\bibitem{Signed} \(\mathbb Z\mathbb Q\) Diophantine Project,
\emph{Signed binomial witnesses force their own positivity and save
five quadratic coordinates}, 23 November 2026.\par
\nolinkurl{undecidability/source_repair_v93/signed_binomial_positivity.tex}.
\bibitem{CRT} \(\mathbb Z\mathbb Q\) Diophantine Project,
\emph{A single integer quotient for the third Pell certificate},
23 November 2026.\par
\nolinkurl{undecidability/pell_crt_v95/one_quotient_pell.tex}.
\end{thebibliography}
\end{document}
